% !TEX root = ../main.tex
\chapter{Topological manifolds}\label{ch:c02}
\begin{paracol}{2}

\begin{sources}
\begin{tabularx}{\linewidth}{@{}l l L@{}}
\toprule
\textbf{Video} & \textbf{Length} & \textbf{Content}\\
\midrule
\seg{2}{1}{Lecture 2: Topological manifolds (Topologische Mannigfaltigkeiten)} & 1:41:06 & Motivation (curves in a set need
continuity); topological spaces (topology, open sets), examples (chaotic, discrete, standard topology on $\R^d$);
continuous maps through preimages and the ``Abitur trick''; the step function as a non-continuous example; open versus
closed; number of topologies on finite sets; topology chosen from the physics (particle tracks); topological manifolds;
charts, chart maps, coordinates, atlases; the circle with the stereographic projection.\\
\bottomrule
\end{tabularx}

\smallskip
\textbf{Transcript} \file{materials/transcripts/C02-V01.txt} (German YouTube auto-captions).
\end{sources}

\begin{reading}[Where this lecture sits in the other notes]
\emph{Geometric Anatomy}: topological spaces (lectures 4 and 5: topologies, open sets, continuity, closed sets,
subspace topology), topological manifolds and atlases (lecture 6).
\emph{Gravity and Light} (TP-1 version): lecture 2 (topology) and lecture 3 (topological manifolds), where the
same definitions are treated with more detail.
\end{reading}

% ==================================================================
\section{Why mathematics first, and why topology}\label{sec:tm-motivation}

Today begins the technical part of the course, which will occupy several hours: the mathematical foundations.
Schuller first wanted to call it ``the structure of spacetime'', but it is rather the \emph{mathematical structure
underlying a concept of space, or of time, or of space and time, or of spacetime} \ts{2}{1}{0:04}.
Physics is not done yet, only mathematics, and it is sensible to separate them: mathematics is something that stays
forever, must make sense in itself and has nothing to do with the world; only because of that status can it be applied
to different situations, by building models of the world. The dirt of the modelling (neglecting air resistance, and so
on) sits in the modelling; the mathematics is clean \ts{2}{1}{0:46}.

Still, a physical motivation from our context \ts{2}{1}{1:29}: the first goal is a mathematical picture of physical
spacetime (do not take the word ``spacetime'' too seriously; what follows holds also for space per se)
\ts{2}{1}{1:29}. At the coarsest, lowest level, spacetime is at least a \emph{set}: intuitively the set of
points \ts{2}{1}{2:03}.

\begin{definition}[New concepts: strategy of structure layers]\label{def:tm-layers}
A set alone does not have enough structure. The key idea is to equip it \emph{successively with more and more
structure}; at each layer of structure there are new things one can do \ts{2}{1}{2:44}.
\end{definition}

For example, in classical physics a point particle should in any case have a \emph{track} (world line, trajectory, curve)
which is \emph{continuous}: it does not jump; particles live infinitely long, we cannot create or annihilate them
(``we could do that, but we do not''); a classical particle does not hop from here to there
\ts{2}{1}{3:04}.

A curve $\gamma$ in a set $M$ is mathematically a \emph{map}
\[
  \gamma\colon \R\to M
\]
that assigns to each parameter value (think of time, but it is just a parameter that runs as you move along the curve) a
point of $M$ \ts{2}{1}{4:29}. Careful: two curves can have the same \emph{trace} but be
different maps. If $\gamma(0)=\tilde\gamma(0)$, $\gamma(1)=\tilde\gamma(1)$ but $\gamma(\tfrac12)\neq\tilde\gamma(\tfrac12)$
(the second curve is traversed faster at first and then slower), the pictures look the same, but the two maps are
different mathematical objects, hence different curves \ts{2}{1}{5:11}.

The big problem: what does \emph{continuous} mean? If the space has no more structure than a set, one cannot even
speak of continuity \ts{2}{1}{6:57}. So one must give ``the naked set without structure'' more
structure \ts{2}{1}{7:42}. School knows the $\varepsilon$--$\delta$ criterion (once secondary-school
material in NRW, ``education republic''), which uses the absolute value on $\R$, in higher dimension a distance or a norm
\ts{2}{1}{8:04}. That is a far too special structure. We want the \emph{weakest} structure
that still allows a clean definition of continuity, and that is a \textbf{topology} \ts{2}{1}{9:07}.
``Weakest structure'' means least prejudice: every structure we add is a prejudice, and the weakest structure is the
smallest prejudice \ts{2}{1}{9:29}.

% ==================================================================
\section{Topological spaces}\label{sec:tm-top}

\begin{definition}[New concepts: topology, open set, topological space]\label{def:tm-topology}
Let $M$ be a set. A collection $\mathcal O\subseteq\mathcal P(M)$ of subsets of $M$ is called a \textbf{topology} on
$M$ (``a topology, not \emph{the} topology: there are probably many possibilities'') if
\ts{2}{1}{10:15}
\begin{enumerate}
  \item $\emptyset\in\mathcal O$ and $M\in\mathcal O$ \ts{2}{1}{11:27};
  \item $U,V\in\mathcal O\ \Rightarrow\ U\cap V\in\mathcal O$ (finite intersections) \ts{2}{1}{12:09};
  \item for any family $(U_\alpha)_{\alpha\in A}$ of elements of $\mathcal O$, with an index set $A$ of any size,
        $\bigcup_{\alpha\in A}U_\alpha\in\mathcal O$ \ts{2}{1}{12:50}.
\end{enumerate}
The elements of $\mathcal O$ are called the \textbf{open sets} (offene Mengen) of $M$ \emph{with respect to}
$\mathcal O$ \ts{2}{1}{15:20}. The pair $(M,\mathcal O)$ is called a \textbf{topological
space} (topologischer Raum) \ts{2}{1}{16:45}.
\end{definition}
\begin{coursenote}[The drawer picture and the index set]
Schuller explains the three conditions with a drawer of subsets: the empty set and the whole set must be in it; if two
sets are in it, so is their intersection; and if you pull out \emph{arbitrarily many} sets (an index set $A$: finitely many,
or countably many, or as many as the real numbers, or worse), their union must also be in the drawer
\ts{2}{1}{11:27}. Countable vs.\ finite vs.\ uncountable is why an arbitrary index set is used: the
rationals can be counted, the reals not \ts{2}{1}{13:31}.
\end{coursenote}
Remarks. Whether a subset is called open depends on the topology: always ask ``open with respect to what?''; the
first topological decision is the choice of $\mathcal O$ (there is a freedom: one can choose this or that
topology as long as the axioms hold) \ts{2}{1}{16:03}. The second is that the object that has structure is
not the set, but the \emph{pair} $(M,\mathcal O)$: ``I give the set structure by choosing something''
\ts{2}{1}{16:24}. The structure is, as far as Schuller knows, also the weakest that allows the definition
of convergence of sequences \ts{2}{1}{14:36}.

\begin{aside}[Why no motivating example first?]
Schuller gives the abstract definition first and the examples afterwards: every motivating example steers the brain
in some direction, perhaps the intended one, but also towards details of the example that have nothing to do with the
abstract definition, and it is very hard to remove such confusions from the head afterwards
\ts{2}{1}{17:26}.
\end{aside}

\begin{example}[Topologies on a three-element set]\label{ex:tm-three}
Let $M=\{1,2,\text{Grillwurst}\}$ (the course joke: ``12 Grillwurst'' is a three-element set with elements
$1$, $2$, ``Bratwurst''; the transcript reads ``12 Kribl Wurz'' [?]; call the third element $w$) \ts{2}{1}{18:07}.
\begin{itemize}
  \item $\mathcal O=\{\emptyset,\{1,2,w\}\}$ is a topology: $\emptyset$ and $M$ are in; the intersection of any two
        members is $\emptyset$ or $M$; the union of any family is $\emptyset$ or $M$ \ts{2}{1}{18:27}.
  \item $\mathcal O=\{\emptyset,\{1\},\{1,2,w\}\}$ is also a topology (the transcript is unclear; this is the natural
        reading) [?] \ts{2}{1}{20:11}.
  \item $\{\emptyset,\{1\},\{2\},\{1,2,w\}\}$ is \textbf{not} a topology: the union $\{1,2\}$ of $\{1\}$ and $\{2\}$ is
        missing, which contradicts condition~3 \ts{2}{1}{20:55}.
  \item A collection without $\emptyset$ (or without $M$) violates condition~1 \ts{2}{1}{21:15}.
\end{itemize}
Already here, four different topologies were written down, but on a set with three elements there are in total $29$
topologies \ts{2}{1}{21:37}.
\end{example}

\begin{definition}[New concepts: chaotic and discrete topology]\label{def:tm-extreme}
Let $M$ be any set. $\mathcal O_{\mathrm{chaotic}}=\{\emptyset,M\}$ is always a topology: the \textbf{coarsest}
(``gröbste'') one, called the \textbf{chaotic topology}; and $\mathcal O_{\mathrm{discrete}}=\mathcal P(M)$, the power set
(the set of all subsets), is also always a topology: the \textbf{finest} one, the \textbf{discrete topology}
\ts{2}{1}{21:57}.
\end{definition}
Both are useless except as extreme examples, which will return once or twice; the useful topologies must be \emph{chosen}
and lie somewhere in the middle \ts{2}{1}{23:02}. More examples on problem sheet~2, exercise~1
\ts{2}{1}{23:23}.

\begin{exercise}[Topologies, problem sheet 2, exercise 1]\label{ex:tm-ex-sheet}
Find further examples and non-examples of topologies on small sets (as above), checking the three axioms.
\end{exercise}

\subsection{The standard topology on $\R^d$}\label{ssec:tm-standard}

Finite sets are ``kindergarten level''; now an extremely important example \ts{2}{1}{24:04}.
$\R^d$ is \emph{defined} as the Cartesian product of $d$ copies of the real line: the set of ordered $d$-tuples
$(a^1,\dots,a^d)$ with each $a^i\in\R$. These tuples have nothing to do with vectors; it is just an ordered
collection of numbers \ts{2}{1}{24:46}. $\R$ has addition and multiplication (it is a
complete field), so there is a great deal of structure on this set that a generic set does not have; this structure is
used to build a topology \ts{2}{1}{25:49}.
\begin{coursenote}[Indices up]
The indices are written \emph{up} ``because I like it, and so do all differential geometers'': sometimes up and sometimes
down, not arbitrarily, and deep wisdom will later be encoded; for now take it as a mere numbering
\ts{2}{1}{26:10}.
\end{coursenote}

\begin{definition}[New concepts: standard topology on $\R^d$]\label{def:tm-standard}
The \textbf{standard topology} $\mathcal O_{\mathrm{std}}$ on $\R^d$ (it exists only on $\R^d$, not on an arbitrary set;
the name is ``standard'' like ``chaotic'' or ``discrete'') is given implicitly: a subset $U\subseteq\R^d$ belongs to
$\mathcal O_{\mathrm{std}}$ if and only if
\[
  \forall p\in U\ \exists r\in\R_{>0}\colon\ B_r(p)\subseteq U,
\]
where the \textbf{soft ball} (``weicher Ball'', an open ball)
\[
  B_r(p):=\Bigl\{q\in\R^d\ \Big|\ \sum_{i=1}^d\bigl(q^i-p^i\bigr)^2<r^2\Bigr\}
\]
has \emph{strict} inequality: points at distance exactly $r$ do not belong to it
\ts{2}{1}{26:54}.
\end{definition}
Notation: $p=(p^1,\dots,p^d)$, often abbreviated $p$; one could also write an arrow or underline it, ``but it has nothing to
do with vectors'' \ts{2}{1}{30:08}. Every topology is a subset of the power set; here the topology is a collection
with uncountably many members, hence not given by an explicit list \ts{2}{1}{27:57}.
Quantifiers: be careful, ``there exists'' and ``for all'' exchange under negation \ts{2}{1}{32:59}.

\begin{example}[Open and not open subsets of $\R^2$]\label{ex:tm-open-r2}
The board is (a piece of) $\R^2$. The set of all points inside a dashed curve, without the dashed boundary, is open in the
standard topology: for each point one can find a radius $r$ with $B_r(p)$ entirely inside \ts{2}{1}{31:54}.
If the boundary (the curve) is included, the set is \emph{not} in the standard topology: at a boundary point, however
small $r>0$ is chosen, the ball contains points outside \ts{2}{1}{33:40}.
\end{example}
That these are exactly the sets of a topology (i.e.\ that this implicit definition satisfies the three axioms) must
still be checked: that is on the next problem sheet \ts{2}{1}{34:24}.
\begin{exercise}[The standard topology is a topology]\label{ex:tm-std-is-top}
Prove that $\mathcal O_{\mathrm{std}}$ of \cref{def:tm-standard} satisfies the three axioms of \cref{def:tm-topology}
(next problem sheet).
\end{exercise}

% ==================================================================
\section{Continuous maps}\label{sec:tm-continuous}

One must choose a topology to have a topological space, and one cannot speak of continuity of maps between two sets unless
both sets carry a topology \ts{2}{1}{35:06}. A map between two bare sets requires no more
structure: it assigns to each element of the domain exactly one element of the target \ts{2}{1}{35:51}.
To call some maps ``continuous'' and others not we need topologies \ts{2}{1}{36:33}.

\begin{definition}[New concepts: continuous map, preimage]\label{def:tm-continuous}
Let $(M,\mathcal O_M)$ and $(N,\mathcal O_N)$ be topological spaces (the topologies may differ because the sets
differ). A map $f\colon M\to N$ is called \textbf{continuous} (stetig) with respect to $\mathcal O_M$ and $\mathcal O_N$
if for every open set $V\in\mathcal O_N$ the \textbf{preimage} (Urbild)
\[
  \operatorname{preim}_f(V)=f^{-1}(V):=\{m\in M\mid f(m)\in V\}
\]
lies in $\mathcal O_M$ \ts{2}{1}{36:54}.
\end{definition}
The preimage is a subset of the domain, defined for \emph{any} subset $V\subseteq N$ (here we use it only for open $V$)
\ts{2}{1}{39:26}. ``Such a simple definition only exists when the concepts are really good.''
A memorable version: \emph{a map is continuous if preimages of open sets are open}, where the first ``open'' refers to
$\mathcal O_N$ (the target) and the second to $\mathcal O_M$ (the domain): the wording hides the fact that these are
different topologies \ts{2}{1}{44:12}. So continuity of a map depends on the map, and on the
topology chosen in the domain \emph{and} in the target \ts{2}{1}{44:57}.

\begin{example}[Men and supermodels]\label{ex:tm-supermodels}
Let $M=\{m_1,m_2,m_3,m_4\}$ be four men, $N=\{\text{Heidi},\text{Claudia},\text{Naomi}\}$ [?: the transcript has
``heidi, klein[?], najdi'', i.e.\ three supermodels; call them $a,b,c$ in what follows] and let $f$ assign to each man
the supermodel he stares at (one at a time, so it is a map): $m_1\mapsto b$, $m_2\mapsto a$ ... [?]. The transcript
leaves the assignment unclear; what is said clearly is how preimages work: the preimage of the whole target set is
all of $M$ (true for every map); the preimage of a one-element set is the set of all men staring at that model; the
preimage of a set containing a model nobody stares at is the empty set (``it is child's play, one only has to draw it
once'') \ts{2}{1}{40:08}.
\end{example}

\begin{example}[The Abitur trick]\label{ex:tm-abitur}
Suppose you must prepare the worst pupil for the oral Abitur: the external examiner gives a very complicated function
$\R\to\R$, e.g.\ $x\mapsto$ (something with sine and cosine), and asks whether it is continuous. Teach the pupil by
heart: ``I equip the domain with the \emph{discrete} topology'': then every map from the domain is continuous (all subsets
are open, so all preimages are open). Cleverer still: tell the examiner ``choose any topology on the target, I take the
discrete topology on the domain'' \ts{2}{1}{45:48}.
Some examiners are fooled twice: they then give a topology on the domain first; if the pupil is fast enough, he takes the
\emph{chaotic} topology on the target, and then every map is continuous (the only open sets are $\emptyset$ and $N$, with
preimages $\emptyset$ and $M$) \ts{2}{1}{47:58}.\par
Both topologies are thus ``a bit useless'': with them every function is continuous, no distinction
continuous/non-continuous remains \ts{2}{1}{48:20}. If you understand this example on the
basis of the definitions of continuity and preimage, everything there is to understand has been understood
\ts{2}{1}{49:02}.
\end{example}

\begin{example}[The soothing example: a step function in the standard topology]\label{ex:tm-step}
Take functions $\R\to\R$ with the standard topology on both $\R=\R^1$ (a soft ball in $\R^1$ is an open interval; soft
balls are open sets, hence also called ``open balls'') \ts{2}{1}{49:23}.
Let $f$ be a function that approaches a positive constant and then \emph{jumps} to a higher value at $x=2$ (the value at
the jump point belongs to the upper level, say), then continues without interesting behaviour \ts{2}{1}{50:46}.
Is it continuous? Intuition says no; now it must be shown with topological means: find \emph{one} open set in the target
whose preimage is not open \ts{2}{1}{51:06}.
Choose an open interval $V$ around the upper value, so small that it does not contain the lower value (in class the open set
is also chosen as a union of two open intervals, one around each level; the result is the same)
\ts{2}{1}{52:29}. Its preimage is the half-line $[2,\infty)$, \emph{including} the point $2$,
since there $f$ takes the upper value; and it is not open: a ball around $2$, however small, contains points left of $2$,
which are mapped to the lower value \ts{2}{1}{53:12}.
(The half-line $(-\infty,2)$ that appears for the lower level is open in the standard topology, \ts{2}{1}{53:54}
which is why one must test the right open set.)
The preimage of an open set is not open, hence $f$ is \emph{not} continuous. ``So it is reassuring: with the standard
topology on both domain and target, the function jumps, and it is not continuous.''
\end{example}
\begin{coursenote}[Reading of the step-function passage]
The auto-captions of this passage (\ts{2}{1}{53:12}--\ts{2}{1}{56:50}) are very garbled (``brown set'', ``yellow set'',
a board picture). What is reconstructed here is the logical content: the preimage of a suitable open set is a half-line
containing its endpoint, hence not open [?: the exact sets drawn on the board may differ].
\end{coursenote}

\begin{definition}[New concepts: closed set]\label{def:tm-closed}
A subset $A\subseteq M$ of a topological space is called \textbf{closed} (abgeschlossen) if its complement
$M\setminus A$ is open \ts{2}{1}{56:50}.
\end{definition}
Remark: ``not open'' does not mean ``closed''. There are sets that are open and not closed, closed and not open, both,
and neither: all combinations occur; the notions are connected through the complement but are not opposites
\ts{2}{1}{57:54}. So the statement ``the set in the example is not open'' does not mean that it is
closed \ts{2}{1}{58:15}.
\begin{exercise}[Open and closed]\label{ex:tm-clopen}
(Finger exercise for lunch.) Show that $\emptyset$ and $M$ are both open and closed in every topological space
\ts{2}{1}{57:32}. Give examples of sets that are open but not closed, closed but not open, neither,
and (apart from $\emptyset,M$) both, for a suitable topology.
\end{exercise}
More ``super-duper'' examples on problem sheet~2, exercise~2: again, continuity depends on the map, and on the topology in
the domain and in the target \ts{2}{1}{58:35}.

% ==================================================================
\section{Back to curves: which topology?}\label{sec:tm-curves}

Back to the motivation \ts{2}{1}{59:16}: a curve (the track of a particle in space or spacetime) is a map
$\gamma\colon\R\to M$, with a parameter that orders the points of the curve one after another and has the taste of time
\ts{2}{1}{59:59}. In mechanics one parametrises the curve by time; but it is not clever to do so from the
beginning, one rather parametrises arbitrarily first and relates this parameter to physical time later; that matters
much more in relativistic mechanics. So Schuller says simply ``the parameter'', and the audience can think what they
like \ts{2}{1}{1:00:20}.

The curve should be at least continuous to describe a classical particle (in quantum physics, one would tell stories about
jumps) \ts{2}{1}{1:01:22}. But we learned today that if $M$ has more than one element, there is more than
one topology; the question to a physicist who demands continuity must be: \emph{with respect to which topology on space or
spacetime}? \ts{2}{1}{1:02:03}

\begin{definition}[Number of topologies on finite sets]\label{def:tm-count}
For a set with $1,2,3,4,5,6,7$ elements, the number of different topologies is
$1,\ 4,\ 29,\ 355,\ 6942,\ 209{,}527,\ 9{,}535{,}241$ [the transcript has ``9 million 535 1241'', ``209.000 ... 527'']
\ts{2}{1}{1:02:45}.
\end{definition}
Imagine how many topologies one can put on $\R^4$ or $\R^3$ (infinitely many) \ts{2}{1}{1:03:49}. There is no simple formula:
the number is bounded by the power set $2^{2^{|M|}}$, but the topologies are subsets of the power set with complicated
conditions (Schuller does not know a closed formula; ``I copied this from Wikipedia'') \ts{2}{1}{1:04:31}.

Where does the wisdom ``spacetime is $\R^4$ with the standard topology'' come from? Are you sure? Schuller says no, and it
is not clear even in classical physics \ts{2}{1}{1:05:13}.

\begin{intuition}[Local physics does not tell the global topology]
You can walk and walk and never reach the end; but imagine flying with a rocket millions of light years away (``has
anyone ever done this? Come to me afterwards''). It is like the earth in earlier times: people walked and walked
(``the Earth is a disc, what else'') until Columbus, and arrived at the other end. So from our local physics and life
experience within a few million light years we do not really know the global structure of space: is it $\R^3$, or is
it perhaps the three-sphere, so that continuing to fly one comes out ``behind'' where one started? The sphere has a
different topology from the plane; on the exercise sheet there is the circle versus a line. The fundamental problem:
``continuous'' sounds great but is without substance as long as we do not say which topology space has, and
\emph{where would I get that}? \ts{2}{1}{1:05:34}.
\end{intuition}

The answer is the same trick as last time (transfer needed) \ts{2}{1}{1:06:56}. In the first law, the
straight lines, defined by force-free particles, gave the geometry. Likewise: to say what a straight line is one needs
more structure (a distance, for example: shortest connection); but we do not have it. The idea of the first law was
to \emph{throw the particles and watch}. So: I want to describe classical particles by continuous curves, so whenever I see
a classical particle I \emph{call its track continuous}, i.e.\ I must choose the topology such that all the tracks of
the potential particles that I observe are continuous. I do not say ``I know the topology, now I know what particles
do'': I let the particle tell me \ts{2}{1}{1:07:17}.
Theorists with ``dogmatic slumber'' like to say that spaces must be clarified first; ``where did you get them, from the
post office?'' The particle tells it \ts{2}{1}{1:08:40}.

\begin{exercise}[Coarsest and finest topology making a map continuous]\label{ex:tm-final-initial}
(Next problem sheet.) For a given topology on one set and a given map $f$ on finite sets, find the \emph{finest}
topology on the other set such that $f$ is still continuous (the chaotic topology on the target would always work, but
here we want the largest); in another exercise, in the opposite direction, find the \emph{coarsest}
topology such that $f$ is just continuous. This is exactly the trick physics needs: do not first specify topologies and
then decide, but the reverse \ts{2}{1}{1:09:01}.
\end{exercise}
Does knowing all potential particle tracks and calling them all continuous determine the topology uniquely? Schuller
thinks it is a hard question; it leads to a real mathematical problem that one could solve before writing the theory down;
``we do not do it now'' \ts{2}{1}{1:10:23}.

% ==================================================================
\section{Topological manifolds}\label{sec:tm-manifolds}

Next: topological manifolds \ts{2}{1}{1:11:04}. Topological spaces were the general notion; topological
manifolds are a very special case: topological spaces that \emph{locally look like a piece of $\R^d$} with the standard
topology \ts{2}{1}{1:12:00}. What does ``locally look like'' mean? Contrary to the
principle ``abstract definition first'', now the example first \ts{2}{1}{1:13:05}.

\begin{example}[The surface of a pretzel]\label{ex:tm-pretzel}
Look at the surface of a pretzel (Brezel, with the salt on it) as a point set, with some topology \ts{2}{1}{1:13:25}.
Go to an arbitrary point $p$ and take an arbitrary open neighbourhood $U$ (a subset of the set that contains $p$ and is
open in the topology) \ts{2}{1}{1:14:08}. ``Looks locally like'' then
means: there is a map $x\colon U\to\R^2$ whose image $x(U)$ is a piece of $\R^2$; we require that $x$ is \emph{bijective}
onto its image and that \emph{both} $x$ and its inverse are continuous
\ts{2}{1}{1:15:10}.
Why demand that the inverse is continuous too? Because a continuous bijection need not have a continuous inverse
(think again of the Abitur example with the target equipped by a different topology) \ts{2}{1}{1:16:14}.
Draw a small heart in the pretzel: under a continuous map it may get kinks and anything can happen, but it is \emph{not torn
apart} \ts{2}{1}{1:16:56}. This holds at every point, so the pretzel surface is a two-dimensional
topological manifold; so is the surface of a pretzel with three holes, and so is the two-sphere $S^2$ (the surface of
a ball; whether the ball has a dent does not matter, it is the same topological space, as we shall see) \ts{2}{1}{1:17:23}.
\end{example}

\begin{example}[A non-manifold: the flying carpet]\label{ex:tm-carpet}
A flying carpet of dimension two, without its edge, is a two-dimensional manifold: at each point, even very near the edge, there
is still a neighbourhood. With the edge included (a carpet with its border) one would have points where there is no such
neighbourhood. And if the carpet has \emph{fringes} (``bad quality, they stick out''), then between the fringes it looks
like a piece of $\R^2$ and on the fringes like a piece of $\R^1$; it cannot decide its dimension, so it is
\emph{not} a topological manifold \ts{2}{1}{1:18:47}.
\end{example}
(The transcript of the carpet example is rough; the meaning is the one given: edge points or fringes destroy the local
resemblance to a fixed $\R^d$.)

\begin{definition}[New concepts: $d$-dimensional topological manifold]\label{def:tm-manifold}
A topological space $(M,\mathcal O)$ is called a \textbf{$d$-dimensional topological manifold}
(Mannigfaltigkeit) if for every point $p\in M$ there exist an open neighbourhood $U\in\mathcal O$ of $p$ and a map
$x\colon U\to x(U)\subseteq\R^d$ such that
\begin{enumerate}
  \item $x$ has an inverse map $x^{-1}\colon x(U)\to U$ (i.e.\ $x$ is invertible onto its image);
  \item $x^{-1}$ is continuous;
  \item $x$ is itself continuous
\end{enumerate}
\ts{2}{1}{1:20:30}.
Here continuity of $x$ and $x^{-1}$ refers to the \emph{subspace topologies} on $U\subseteq M$ and on
$x(U)\subseteq\R^d$ (the latter from the standard topology of $\R^d$) \ts{2}{1}{1:25:17}.
\end{definition}
There must be one open neighbourhood \emph{around every point} (maybe several, but at least one), and the language must be
precise (``so that'' three things hold) \ts{2}{1}{1:22:42}.
\begin{coursenote}[Subspace topology]
A subset of a topological space inherits a topology from it; the topology on $U\subseteq M$ and on $x(U)\subseteq\R^d$
is the \emph{subset (subspace) topology}, whose definition and basic properties are on the next problem sheet
(exercise~3, ``inheritance of a topology to a subset'') \ts{2}{1}{1:25:17}. The last
exercise (``the circle as a topological manifold'') also contains charts and atlases \ts{2}{1}{1:26:11}.
\end{coursenote}
\begin{exercise}[Subspace topology and the circle]\label{ex:tm-subspace}
Problem sheet 3, exercise~3: the topology inherited by a subset; last exercise: the circle as a topological manifold.
\end{exercise}

% ==================================================================
\section{Charts and atlases}\label{sec:tm-charts}

The ``babble'' about charts and atlases of the first lecture is now turned into definitions
\ts{2}{1}{1:26:32}.

\begin{definition}[New concepts: chart, chart domain, chart map, coordinates]\label{def:tm-chart}
Let $(M,\mathcal O)$ be a $d$-dimensional topological manifold (one must always say which dimension). A pair
$(U,x)$ is called a \textbf{chart} (Karte) of $M$ if \ts{2}{1}{1:26:55}
\begin{enumerate}
  \item $U\in\mathcal O$ is open;
  \item $x\colon U\to x(U)\subseteq\R^d$ is continuous;
  \item $x^{-1}\colon x(U)\to U$ exists and is continuous
\end{enumerate}
\ts{2}{1}{1:28:25}. The set $U\subseteq M$ is the \textbf{chart domain} (Kartengebiet), the map $x$
the \textbf{chart map}, $x(U)\subseteq\R^d$ the \textbf{chart image}; the number $x^i(p)$ (the $i$-th entry of the
$d$-tuple $x(p)=(x^1(p),\dots,x^d(p))$) is the \textbf{$i$-th coordinate} of $p$ with respect to the chart, and
$x^i\colon U\to\R$ the \textbf{$i$-th coordinate function} \ts{2}{1}{1:29:10}.
\end{definition}
The definition is first made, and then interpreted: the chart domain is the piece in the ``real world'' (the Amazon,
if you drive through it with a road atlas: the Amazon is the chart domain, not what is printed; the real world has
crocodiles); $x(U)$ is the chart image, the nicely printed Amazon in the atlas \ts{2}{1}{1:29:10}.
The $x^i$ assign to a point of the real world numbers \ts{2}{1}{1:31:00}.

\begin{intuition}[Coordinates depend on the chart; physics must not]
The Eiffel tower has coordinates on page 2\,100 of the atlas, and different ones on another map of the same atlas
\ts{2}{1}{1:32:04}. The Eiffel tower does not care about the coordinates one gives it:
imagine it did, then the \emph{physics} of a chart-dependent world would be telekinesis. One could say ``the object does
not move in this chart (the coordinates do not change)''; another person in another chart says it moves; but the object is
the same, and if the object cared about what each of us thinks, we could move it by thinking. So we must formulate
physics \emph{independently of the chart}: if the laws depended on the charts (which are figments, each publisher can make
another atlas), we would have a world with telekinesis (``then you would win a Nobel Prize'': if it were
found, it would have to be shown by experiment). Since we do not want a physics with telekinesis, \emph{everything must be
made chart-independent}: that is the point of chart-independence \ts{2}{1}{1:32:47}.
\end{intuition}
The $i$-th coordinate of $p$ with respect to the chart is $x^i(p)$ (``also the $i$-th coordinate function'', terminology
to give intuition) \ts{2}{1}{1:34:49}.

\begin{definition}[New concepts: atlas]\label{def:tm-atlas}
A collection $\mathcal A$ of charts is called an \textbf{atlas} (Atlas) if the union of the chart domains of all charts
$(U,x)\in\mathcal A$ is the whole manifold: $\bigcup_{(U,x)\in\mathcal A}U=M$ \ts{2}{1}{1:35:12}.
\end{definition}
The intuition: enough chart domains that their union covers the manifold; if a chart is missing, i.e.\ you cannot find
your region in the atlas, it is not an atlas \ts{2}{1}{1:36:17}.

\begin{example}[The circle $S^1$ with stereographic projection]\label{ex:tm-circle}
Take as manifold the circle $S^1$ and cover it with charts satisfying all conditions; there are arbitrarily many ways to
do so; one is the \emph{stereographic projection} (``a frightening name, but childishly easy'')
\ts{2}{1}{1:36:39}.\par
Draw a line below the circle, the target $\R^1$ of a one-dimensional chart. To each point $p$ of the circle assign
a number as follows: start at the \emph{south pole}, draw a straight line through $p$, and see where it hits the line; the
point where the vertical through the centre hits it is the zero; the number there is $x^1(p)$ (for another point $q$ one
gets $x^1(q)$ and so on) \ts{2}{1}{1:37:22}.
Which point of the manifold do we not reach? The south pole itself: the line from the south pole to the south pole does
not hit the target line. The chart domain is thus the whole circle \emph{minus} one point, the south pole
\ts{2}{1}{1:38:46}.\par
No problem: do the same trick again with a second line and start from the \emph{north pole}. Then the north pole is
missing but the south pole is covered (``you see penguins''), and a given point like the Eiffel tower has different
coordinates in the second chart than in the first. The gelbe chart shows no polar bears and the blue one no penguins,
``you know how it goes'' \ts{2}{1}{1:39:28}. The two charts together
form an atlas of the circle. With the circle embedded in $\R^2$ one can write the formulas explicitly: the coordinates of
a point of the circle in the two charts; this is the last exercise on the current problem sheet
\ts{2}{1}{1:40:29}.
\end{example}
\begin{exercise}[Stereographic charts of the circle]\label{ex:tm-stereo}
For the unit circle in $\R^2$ derive the two chart maps $x_{\pm}$ (projection from the south and the north pole) and
their inverses, and check that they satisfy \cref{def:tm-chart}; show the two chart domains form an atlas
(problem sheet, last exercise).
\end{exercise}

% ==================================================================
\section*{Key concepts}
\begin{keyconcepts}
\textbf{Retrieval questions.}
\begin{enumerate}
  \item Why is a bare set not enough to speak of continuity of a curve? What is the weakest structure that is?
  \item State the three axioms of a topology. Which are the coarsest and the finest topologies?
  \item Define the standard topology on $\R^d$. Which subsets of $\R^2$ (with and without boundary) belong to it?
  \item Define a continuous map between topological spaces. Why is every map continuous if the domain is discrete or the
        target chaotic?
  \item Show that a step function $\R\to\R$ is not continuous in the standard topologies.
  \item What does closed mean? Give a set open and not closed, and one neither.
  \item Why should the topology be \emph{chosen} from the particle tracks rather than assumed?
  \item Define a $d$-dimensional topological manifold, a chart, a chart map, coordinates, an atlas. Why is the flying
        carpet with fringes not a manifold?
  \item Why must physics be chart-independent? Give the two charts of the circle.
\end{enumerate}
\textbf{Concepts.}
\begin{itemize}
  \item Layers of structure: set $\to$ topology (\cref{def:tm-layers}); topology, open sets, topological space
        (\cref{def:tm-topology}); chaotic, discrete (\cref{def:tm-extreme}); standard topology (\cref{def:tm-standard}).
  \item Continuity via preimages (\cref{def:tm-continuous}); closed sets (\cref{def:tm-closed}); number of
        topologies on finite sets (\cref{def:tm-count}).
  \item Topological manifold (\cref{def:tm-manifold}); chart, chart domain, chart map, coordinates
        (\cref{def:tm-chart}); atlas (\cref{def:tm-atlas}); the circle (\cref{ex:tm-circle}).
\end{itemize}
\end{keyconcepts}

\flushside
\end{paracol}
